\section*{Глава 2. Проектирование веб-приложения}
\addcontentsline{toc}{section}{Глава 2. Проектирование веб-приложения}
\stepcounter{section}


\subsection{Анализ целевой аудитории и формирование требований}

Анализ целевой аудитории позволяет выделить два основных сегмента пользователей. Первый сегмент — это розничные покупатели, 
заинтересованные в приобретении цифровой техники. Их главная потребность заключается в удобном поиске по специфическим техническим 
характеристикам и возможности сравнения предложений от разных поставщиков. Второй сегмент — представители малого и среднего 
бизнеса, нуждающиеся в доступном канале продажи своей продукции с минимальными техническими барьерами для входа.

Исходя из потребностей целевой аудитории, в проектируемой системе выделяются три ключевые ролевые модели с изолированными правами доступа:

\begin{itemize}
	\item \textbf{Покупатель} --- авторизованный пользователь, имеющий доступ к витрине товаров. Нуждается в инструментах фильтрации, 
			управления корзиной заказов и отслеживания статусов текущих покупок через личный кабинет.
	\item \textbf{Продавец} --- коммерческий пользователь, управляющий своими карточками товаров. Основные потребности продавца: публикация товаров 
			с загрузкой медиафайлов, управление остатками, обработка поступивших заказов и получение агрегированной статистики продаж.
	\item \textbf{Администратор} --- сотрудник платформы, обладающий полными правами. Отвечает за модерацию контента, блокировку нарушителей 
			и управление глобальным справочником категорий техники.
\end{itemize}

На основе выделенных ролей был сформирован перечень различных требований к разрабатываемому программному продукту.

\textbf{Функциональные требования}:
\begin{enumerate}
    \item Регистрация, аутентификация и авторизация пользователей с разграничением прав доступа на основе ролей.
    \item Реализация CRUD-интерфейсов для управления профилями, каталогом товаров и категориями.
    \item Наличие логики управления корзиной покупок и оформления заказа.
    \item Маршрутизация заказов: продавец должен иметь доступ только к тем заказам, которые содержат добавленные им товары.
    \item Возможность загрузки и хранения медиафайлов.
    \item Выполнение агрегации данных на стороне сервера для формирования сводной аналитической отчетности продавца (подсчет выручки, количества заказов).
\end{enumerate}

\textbf{Нефункциональные требования}:
\begin{enumerate}
    \item Производительность: клиентская часть должна функционировать как одностраничное приложение (SPA), обеспечивая фильтрацию каталога без перезагрузки страницы.
    \item Безопасность: пароли пользователей должны храниться в базе данных в виде криптографических хешей; передача сессионных данных должна осуществляться через подписанные JWT-токены.
    \item Масштабируемость: структура базы данных должна быть нормализована для предотвращения аномалий при добавлении новых продавцов или характеристик товара.
	\item Эргономичность: понятный интерфейс, удобная навигация и логичность расположения элементов в приложении.
\end{enumerate}

Таким образом, проведенный анализ целевой аудитории и составление требований формируют четкий вектор для последующего моделирования бизнес-процессов и 
проектирования физической структуры реляционной базы данных.

\subsection{Описание функциональности платформы}

Для формализации функциональных требований и проектирования бизнес-логики разрабатываемой веб-платформы применена методология пользовательских историй. Данный подход 
позволяет описать функциональность программного продукта с позиции конечного потребителя. Рассмотрим ключевые сценарии взаимодействия для каждой роли в системе.

Сценарии для базовой роли \textbf{«Пользователь системы»}:
\begin{itemize}
    \item \textit{Как новый пользователь}, я хочу зарегистрироваться в системе, чтобы получить доступ к функционалу покупателя или продавца.
    \item \textit{Как зарегистрированный пользователь}, я хочу авторизоваться в системе по своим учетным данным, чтобы безопасно управлять своим профилем.
\end{itemize}

Сценарии для роли \textbf{«Покупатель»}:
\begin{itemize}
    \item \textit{Как покупатель}, я хочу просматривать витрину и использовать фильтры, чтобы быстро находить нужную мне электронику.
    \item \textit{Как покупатель}, я хочу переходить в карточку товара и добавлять его в корзину, чтобы формировать список покупок.
    \item \textit{Как покупатель}, я хочу оформлять заказ из корзины и отслеживать его статус в личном кабинете, чтобы контролировать процесс доставки.
\end{itemize}

Сценарии для роли \textbf{«Продавец»}:
\begin{itemize}
    \item \textit{Как продавец}, я хочу иметь доступ к единой панели управления, чтобы централизованно вести коммерческую деятельность.
    \item \textit{Как продавец}, я хочу управлять товарами и загружать к ним медиафайлы, чтобы качественно оформлять витрину.
    \item \textit{Как продавец}, я хочу изменять статусы заказов и просматривать статистику продаж с возможностью выгрузки отчета в формате CSV, чтобы анализировать выручку.
\end{itemize}

Сценарии для роли \textbf{«Администратор»}:
\begin{itemize}
    \item \textit{Как администратор}, я хочу управлять категориями техники, чтобы поддерживать структурированность каталога.
    \item \textit{Как администратор}, я хочу модерировать учетные записи пользователей, чтобы пресекать нарушения правил платформы.
\end{itemize}

На основе сформированных функциональных требований и пользовательских сценариев была разработана UML-диаграмма вариантов использования (рисунок ~\ref{fig:usecase}). Данная модель 
визуализирует границы маркетплейса и детализирует архитектуру взаимодействия ролей с использованием механизмов наследования и расширенной нотации связей.

\begin{figure}[H]
    \centering
    \includegraphics[width=1.0\textwidth]{usecase.png}
    \caption{Диаграмма вариантов использования (Use Case)}
    \label{fig:usecase}
\end{figure}

В системе реализована иерархия акторов: базовая абстракция «Пользователь системы» берет на себя общие прецеденты авторизации и регистрации, а остальные роли наследуют эти 
права. Бизнес-логика изобилует сложными зависимостями. Например, панель управления у продавцов включает в себя модули управления заказами, товарами и просмотра 
статистики. В свою очередь, создание карточки товара обязательно требует загрузки медиафайлов. Отношения расширения применены для дополнительного функционала, который вызывается 
опционально: выгрузка CSV-отчета из статистики или отслеживание заказа после его оформления. Разработанная поведенческая модель послужит надежным фундаментом для следующего 
этапа — проектирования физической структуры реляционной базы данных.

\subsection{Проектирование модели данных}

При проектировании инфологической модели были выделены следующие ключевые сущности:
\begin{itemize}
    \item \texttt{roles} и \texttt{users} — таблицы для реализации системы авторизации и управления правами доступа.
    \item \texttt{categories} и \texttt{products} — таблицы для формирования структуры каталога и хранения карточек товаров. Каждая карточка товара жестко привязана к продавцу и категории.
    \item \texttt{orders} и \texttt{order\_items} — таблицы для реализации бизнес-логики покупок. 
\end{itemize}

Особое внимание при проектировании было уделено архитектурной проблеме заказов: в классическом интернет-магазине статус присваивается всему заказу целиком, однако в маркетплейсе один заказ может содержать товары от разных продавцов. Для решения этой проблемы статус заказа был вынесен на уровень позиций заказа. Это позволяет каждому продавцу изменять статус доставки только своего товара, не влияя на общую корзину покупателя.

Логическая структура базы данных со связями типа «один-ко-многим» представлена на ER-диаграмме (рисунок~\ref{fig:erd}).

\begin{figure}[H]
    \centering
    \includegraphics[width=0.8\textwidth]{erd.png}
    \caption{ER-диаграмма базы данных}
    \label{fig:erd}
\end{figure}

Для программной реализации спроектированная модель была нормализована до 3-й нормальной формы. Ниже (таблицы~\ref{tab:db_users}, \ref{tab:db_products}, \ref{tab:db_orders}, \ref{tab:db_order_items}) представлены физические схемы основных таблиц базы данных.

{
\small
\linespread{1.0}\selectfont
\renewcommand{\arraystretch}{1.1}

% --- Таблица USERS ---
\begin{longtable}{|p{3.2cm}|p{3cm}|p{3.5cm}|p{5.3cm}|}
	\caption{Структура таблицы пользователей (\texttt{users})} \label{tab:db_users} \\
	\hline
	\textbf{Поле} & \textbf{Тип данных} & \textbf{Ограничение} & \textbf{Описание} \\ \hline
	\endfirsthead

	\multicolumn{4}{r}{\textit{Продолжение таблицы \ref{tab:db_users}}} \\
	\hline
	\textbf{Поле} & \textbf{Тип данных} & \textbf{Ограничение} & \textbf{Описание} \\ \hline
	\endhead

	\hline
	\endfoot

	\hline
	\endlastfoot

	id & UUID & PRIMARY KEY & Уникальный идентификатор \\ \hline
	role\_id & INT & FOREIGN KEY & Ссылка на роль пользователя \\ \hline
	login & VARCHAR(50) & UNIQUE, NOT NULL & Логин для входа \\ \hline
	email & VARCHAR(255) & UNIQUE, NOT NULL & Электронная почта \\ \hline
	password\_hash & VARCHAR(255) & NOT NULL & Хеш пароля (bcrypt) \\ \hline
	first\_name & VARCHAR(100) & NOT NULL & Имя \\ \hline
	last\_name & VARCHAR(100) & NOT NULL & Фамилия \\ \hline
	middle\_name & VARCHAR(100) & NULL & Отчество (необязательно) \\ \hline
	phone\_number & VARCHAR(20) & NULL & Телефон для связи \\ \hline
	created\_at & TIMESTAMP & DEFAULT NOW() & Дата регистрации \\ \hline
\end{longtable}

% --- Таблица PRODUCTS ---
\begin{longtable}{|p{3.2cm}|p{3cm}|p{3.5cm}|p{5.3cm}|}
	\caption{Структура таблицы товаров (\texttt{products})} \label{tab:db_products} \\
	\hline
	\textbf{Поле} & \textbf{Тип данных} & \textbf{Ограничение} & \textbf{Описание} \\ \hline
	\endfirsthead

	\multicolumn{4}{r}{\textit{Продолжение таблицы \ref{tab:db_products}}} \\
	\hline
	\textbf{Поле} & \textbf{Тип данных} & \textbf{Ограничение} & \textbf{Описание} \\ \hline
	\endhead

	\hline
	\endfoot

	\hline
	\endlastfoot

	id & UUID & PRIMARY KEY & Уникальный идентификатор \\ \hline
	vendor\_id & UUID & FOREIGN KEY & Ссылка на продавца \\ \hline
	category\_id & INT & FOREIGN KEY & Ссылка на категорию \\ \hline
	title & VARCHAR(255) & NOT NULL & Наименование товара \\ \hline
	description & TEXT & NULL & Подробное описание \\ \hline
	price & DECIMAL(10,2) & CHECK > 0 & Стоимость товара \\ \hline
	stock\_quantity & INT & DEFAULT 0 & Остаток на складе \\ \hline
	image\_url & VARCHAR(255) & NULL & Путь к файлу изображения \\ \hline
	attributes & JSONB & NOT NULL & Динамические характеристики \\ \hline
\end{longtable}

% --- Таблица ORDERS ---
\begin{longtable}{|p{3.2cm}|p{3cm}|p{3.5cm}|p{5.3cm}|}
	\caption{Структура таблицы заказов (\texttt{orders})} \label{tab:db_orders} \\
	\hline
	\textbf{Поле} & \textbf{Тип данных} & \textbf{Ограничение} & \textbf{Описание} \\ \hline
	\endfirsthead

	\multicolumn{4}{r}{\textit{Продолжение таблицы \ref{tab:db_orders}}} \\
	\hline
	\textbf{Поле} & \textbf{Тип данных} & \textbf{Ограничение} & \textbf{Описание} \\ \hline
	\endhead

	\hline
	\endfoot

	\hline
	\endlastfoot

	id & UUID & PRIMARY KEY & Уникальный идентификатор \\ \hline
	customer\_id & UUID & FOREIGN KEY & Ссылка на покупателя \\ \hline
	delivery\_address & TEXT & NOT NULL & Полный адрес доставки \\ \hline
	total\_amount & DECIMAL(10,2) & NOT NULL & Итоговая сумма заказа \\ \hline
	created\_at & TIMESTAMP & DEFAULT NOW() & Дата и время оформления \\ \hline
\end{longtable}

% --- Таблица ORDER_ITEMS ---
\begin{longtable}{|p{3.2cm}|p{3cm}|p{3.5cm}|p{5.3cm}|}
	\caption{Структура таблицы позиций заказа (\texttt{order\_items})} \label{tab:db_order_items} \\
	\hline
	\textbf{Поле} & \textbf{Тип данных} & \textbf{Ограничение} & \textbf{Описание} \\ \hline
	\endfirsthead

	\multicolumn{4}{r}{\textit{Продолжение таблицы \ref{tab:db_order_items}}} \\
	\hline
	\textbf{Поле} & \textbf{Тип данных} & \textbf{Ограничение} & \textbf{Описание} \\ \hline
	\endhead

	\hline
	\endfoot

	\hline
	\endlastfoot

	id & UUID & PRIMARY KEY & Уникальный идентификатор \\ \hline
	order\_id & UUID & FOREIGN KEY & Ссылка на общий чек (заказ) \\ \hline
	product\_id & UUID & FOREIGN KEY & Ссылка на купленный товар \\ \hline
	quantity & INT & CHECK > 0 & Количество единиц товара \\ \hline
	price\_at\_purchase & DECIMAL(10,2) & NOT NULL & Цена на момент покупки \\ \hline
	status & ENUM & NOT NULL & Статус отправки продавцом \\ \hline
\end{longtable}
}

Спроектированная реляционная модель исключает дублирование данных и полностью готова к интеграции с серверной частью приложения через ORM SQLAlchemy.

\subsection{Разработка макетов страниц}

Для визуализации пользовательского интерфейса и проектирования путей взаимодействия были разработаны низкодетализированные структурные макеты основных 
страниц (Wireframes). Данный подход позволяет сфокусироваться на эргономике, расположении элементов управления и информационной архитектуре до начала 
этапа программной верстки. Проектирование макетов осуществлялось с учетом компонентного подхода, который будет реализован в дальнейшем.

Главная страница является точкой входа для покупателей. В верхней части расположена глобальная панель навигации с модулем глобального поиска, кнопками 
авторизации и доступа к корзине. Ниже расположен рекламный блок с призывом к действию (CTA). Ключевым элементом интерфейса является панель фильтрации 
и блок категорий, позволяющие покупателю отсеивать технику по заданным критериям. Основное пространство занимает модульная сетка популярных товаров. Макет 
главной страницы представлен на рисунке~\ref{fig:mock_main}.

\begin{figure}[H]
    \centering
    \includegraphics[width=0.8\textwidth]{mock_main.png}
    \caption{Макет главной страницы}
    \label{fig:mock_main}
\end{figure}

Карточка товара разработана с акцентом на детальное представление сложной электроники. Макет (рисунок~\ref{fig:mock_itempage}) включает крупную область для просмотра 
фотографий товара. В информационной панели, помимо цены и названия, предусмотрен вывод имени продавца и динамических селекторов характеристик. Ниже располагаются 
блоки с текстовым описанием, техническими спецификациями и рекомендательной системой товаров из той же категории.

\begin{figure}[H]
    \centering
    \includegraphics[width=0.8\textwidth]{mock_itempage.png}
    \caption{Макет карточки товара}
    \label{fig:mock_itempage}
\end{figure}

Интерфейс корзины (рисунок~\ref{fig:mock_shopcart}) спроектирован для удобной агрегации товаров от различных продавцов. Пользователь имеет возможность управлять 
количеством позиций и удалять неактуальные товары. Правая боковая панель зафиксирована и отображает информацию о заказе с кнопкой перехода к оформлению заказа.

\begin{figure}[H]
    \centering
    \includegraphics[width=0.8\textwidth]{mock_shopcart.png}
    \caption{Макет корзины}
    \label{fig:mock_shopcart}
\end{figure}

Для продавцов спроектирован изолированный интерфейс панели управления (рисунок~\ref{fig:mock_dash_vendor}). Макет имеет классическую для B2B-систем структуру, соответствующую 
современным паттернам проектирования коммерческих интерфейсов\cite{ux_b2b_design}: левое боковое меню для навигации между разделами и широкую рабочую область. В ней визуализированы 
графики агрегированной статистики продаж с функцией экспорта данных в формате CSV. Ниже расположены списки для выполнения CRUD-операций над карточками товаров и управления 
статусами поступивших заказов.

\begin{figure}[H]
    \centering
    \includegraphics[width=0.8\textwidth]{mock_dash_vendor.png}
    \caption{Макет панели управления продавца}
    \label{fig:mock_dash_vendor}
\end{figure}

Таким образом, завершен этап архитектурного планирования, и подготовлена исчерпывающая база для программной реализации серверной и клиентской частей маркетплейса в следующей главе.